% Appendix A

\chapter{Installation and Execution Instructions}
\label{AppendixA}
\lhead{Appendix A. \emph{Installation and Execution Instructions}}

\section{Submitted Folders and Files}
The submitted project package should include the following folders and files:
\begin{itemize}
    \item \texttt{crates/}: Rust workspace crates for the host daemon, guest
    runtime, shared protocol, security, memory, tools, heartbeat, and channels.
    \item \texttt{configs/}: runtime configuration files for IronClaw and the
    guest runtime.
    \item \texttt{scripts/}: build and smoke-test scripts for local and
    Firecracker-backed execution.
    \item \texttt{rootfs/}: guest root filesystem skeleton files.
    \item \texttt{kernels/firecracker/}: kernel images used by the Firecracker
    runtime.
    \item \texttt{docs/}: supporting architecture and setup documentation.
    \item \texttt{experiments/}: H1 TypeScript harness, deterministic fixtures,
    raw JSON results, and rendered H1/H2 result reports.
    \item \texttt{Report Thesis Suryavir Kapur/}: thesis report source and
    generated PDF.
\end{itemize}

\section{Prerequisites}
The development machine should have Rust installed through \texttt{rustup}, the
toolchain specified in \texttt{rust-toolchain.toml}, \texttt{cargo},
\texttt{make}, a Linux environment capable of running Firecracker, and LaTeX
tools for building the report. Firecracker execution may require KVM support and
appropriate permissions on the host.

\section{Compilation Steps}
From the repository root, build the Rust workspace using:
\begin{verbatim}
cargo build --workspace
\end{verbatim}

Run the test suite using:
\begin{verbatim}
cargo test --workspace
\end{verbatim}

Check every feature and target, including the Firecracker evaluation example:
\begin{verbatim}
cargo check --workspace --all-targets --all-features
\end{verbatim}

Build the thesis PDF from the report directory using:
\begin{verbatim}
cd "Report Thesis Suryavir Kapur"
make
\end{verbatim}

\section{Execution Steps}
Local smoke tests can be executed using the scripts under \texttt{scripts/}.
Typical commands include:
\begin{verbatim}
./scripts/smoke-local-guest.sh
./scripts/smoke-firecracker.sh
./scripts/smoke-firecracker-ws.sh
./scripts/smoke-guest-tools-ws.sh
\end{verbatim}

Configuration files under \texttt{configs/} select the runtime behavior,
permissions, and guest execution mode. Before running Firecracker tests, verify
that the configured kernel path, root filesystem path, vsock settings, and file
permissions match the local machine.

\section{Evaluation Reproduction}
Reproduce H1 from the \texttt{experiments/} directory:
\begin{verbatim}
export DEEPSEEK_API_KEY=...
npm install
npm run experiment
\end{verbatim}
This regenerates \texttt{RESULTS.md} and \texttt{results.json}. Individual
conditions can be selected, for example \texttt{npm run experiment -{}- E5:B}.

Reproduce H2 from the repository root on a Linux host with KVM, Firecracker,
the documented guest kernel, and the immutable Ubuntu root filesystem:
\begin{verbatim}
cargo run -p common --features firecracker \
  --example h2_security_performance
\end{verbatim}
This writes \texttt{experiments/H2\_RESULTS.md} and
\texttt{experiments/h2-results.json}.

\section{Configuration Details}
The main configuration files are:
\begin{itemize}
    \item \texttt{configs/ironclawd.firecracker.toml}: Firecracker-backed host
    daemon configuration.
    \item \texttt{configs/ironclawd.guest-tools.toml}: guest-tools execution
    mode.
    \item \texttt{configs/ironclawd.guest-autonomous.toml}: autonomous guest
    execution mode.
    \item \texttt{configs/irowclaw.allow-bash.toml}: guest runtime configuration
    that permits shell execution for controlled test scenarios.
\end{itemize}

\section{Reproducibility Checklist}
The final project package should make it possible to reproduce the implementation
and evaluation environment. Before submission, verify that:
\begin{itemize}
    \item the Rust workspace builds with \texttt{cargo build --workspace};
    \item the available automated tests run with \texttt{cargo test --workspace};
    \item configuration files under \texttt{configs/} match the local execution
    mode;
    \item Firecracker kernel and root filesystem paths are documented;
    \item smoke-test commands and expected outputs are recorded;
    \item H1 and H2 result JSON files can be regenerated with the commands above;
    \item generated report and presentation PDFs are included.
\end{itemize}
